\par\noindent\textbf{(1)} Jacobi–Anger恒等式exp[iβsin(Ωt)]=Σ\_\{n=−∞\}∞J\_n(β)exp(inΩt)。故输出频率为ωn=ω0+nΩ，复振幅为E0J\_n(β)，n为整数；相位调制不是仅产生一个位于ω0+Ω的边带。
\par\noindent\textbf{(2)} 边带功率分数P\_n/P\_total=J\_n(β)²，J\_\{−n\}(β)=(−1)\textasciicircum{}nJ\_n(β)，故上下n阶功率相等而相位可能相差π。由于每个时间的|E(t)|²=|E0|²为常数，理想相位调制没有改变平均总光功率。
\par\noindent\textbf{(3)} Bessel完备恒等式Σ\_\{n∈Z\}J\_n(β)²=1保证谱总功率守恒；对称性给ΣnJ\_n²=0，而二阶谱矩Σn²J\_n(β)²=β²/2，边带角频率标准差为Ω|β|/√2，与时域瞬时频率ω\_inst=ω0+βΩcosΩt的方差一致。
\par\noindent\textbf{(4)} 当|β|≪1，J0≈1−β²/4，J±1≈±β/2，故首对边带各占β²/4；大β时显著边带延伸到|n|约β，带宽量级2βΩ。若有色散传播或调制器振幅残差，频谱对称与功率分配会改变。
